\documentclass[slidestop,compress,mathserif]{beamer}
%\documentclass[slidestop,compress,mathserif,handout]{beamer}

%\documentclass[xcolor=dvipsnames,handout]{beamer}
%\documentclass[xcolor=dvipsnames]{beamer}

%\documentclass[handout]{beamer}

%%% To get rid of solutions on handouts:
\newcommand{\soln}[1]{\textit{\textcolor{darkGray}{#1}}}				% For slides
%\newcommand{\soln}[1]{ }	% For handouts

% to get pausing to work properly on slides
\newcommand{\hide}[1]{#1}	% For slides
%\newcommand{\hide}[1]{ }	% For handouts


%\usepackage{multicol}
\usepackage{amsfonts}
%\usepackage[pdftex,dvipsnames]{color}
\usepackage{graphicx}
\usepackage{subfigure}
%\usepackage{picinpar}
\usepackage{pifont}
\usepackage{pgf,pgfarrows,pgfnodes}
%\usepackage{wasysym,manfnt,phaistos,empheq}
\usepackage[english]{babel}
\usepackage{pgfpages}
\usepackage{natbib}
\usepackage{hyperref}
\usepackage{multimedia}
%\usepackage{amsfonts,amstext,amssymb,amsbsy,amsopn,amsthm,eucal,latexsym,mathrsfs}
\usepackage{amsmath,amsfonts,amstext,amssymb,amsbsy,amsopn,amsthm,eucal,latexsym,mathrsfs}
\usepackage{ulem}
\usepackage{setspace}
\usepackage{array}
%\usepackage{rotating}
\usepackage{multirow}
\usepackage{verbatim}
\usepackage{multicol}

\setbeamertemplate{navigation symbols}{}

%\usepackage{tikz}
%\usetikzlibrary{arrows,shapes,trees,backgrounds}


%\setbeameroption{show notes on second screen}
%\setbeameroption{show notes}
%\setbeameroption{show only notes}

\definecolor{links}{HTML}{2A1B81}
\hypersetup{colorlinks,linkcolor=,urlcolor=links}

\newtheorem*{principle}{Inscrutibility Principle}
\newtheorem*{punchline}{Punch Line}
\newtheorem{defn}{Definition}

\definecolor{Scarlet}{RGB}{140,17,17}
\definecolor{VassarRed}{RGB}{128,0,0}

% "dinglist" environment
  \renewenvironment{dinglist}[2][blue]
  {\begin{list}{\textcolor{blue}{\ding{#2}}}{}}{\end{list}}
  % Symbol definitions for these lists
  \newcommand{\DingListSymbolA}{43}
  \newcommand{\DingListSymbolB}{243}
  \newcommand{\DingListSymbolC}{224}
  \newcommand{\DingListSymbolD}{219}
  \newcommand{\DingListSymbolCheck}{52}
  \newcommand{\DingListSymbolCross}{56}


  \newenvironment{ballotenv}
{\only{%
\setbeamertemplate{itemize item}{\ding{45}}%
\setbeamertemplate{itemize subitem}{\ding{46}}%
\setbeamertemplate{itemize subsubitem}{\ding{46}}}} {}
\setbeamertemplate{itemize item}{\ding{49}}
\setbeamertemplate{itemize subitem}{\ding{47}}
\setbeamertemplate{itemize subsubitem}{\ding{47}}


%User defined colors: See colors section
\xdefinecolor{oiBlue}{rgb}{0.15, 0.35, 0.55}
\xdefinecolor{gray}{rgb}{0.5, 0.5, 0.5}
\xdefinecolor{darkGray}{rgb}{0.3, 0.3, 0.3}
\xdefinecolor{darkerGray}{rgb}{0.2, 0.2, 0.2}
\xdefinecolor{rubineRed}{rgb}{0.89,0,0.30}
\xdefinecolor{linkCol}{rgb}{0.11,0.49,0.95}	
\xdefinecolor{irishGreen}{rgb}{0,0.60,0}	
\xdefinecolor{darkturquoise}{rgb}{0.44, 0.58, 0.86}
\definecolor{lightGreen}{rgb}{0.533,0.765,0.42}
\xdefinecolor{Regalia}{HTML}{522D80}
\xdefinecolor{ClemsonOrange}{HTML}{EA6A20}

\definecolor{duke@LightGrey}{RGB}{200,200,200}\definecolor{DarkGreen}{RGB}{0,100,0}
\definecolor{Oranges}{RGB}{255,127,0}
\definecolor{LightGray}{RGB}{211,211,211}

%\setbeamertemplate{footline}{%
%  \raisebox{5pt}{\makebox[\paperwidth]{\hfill\makebox[10pt]{\scriptsize\insertframenumber}}}}

\setbeamercolor{equation background}{fg=black,bg=duke@LightGrey}
  % Boxed equation
  \newcommand{\eqbox}[2][0.6]{%
  \centerline{
  \begin{beamerboxesrounded}[lower=equation background,width=#1\hsize,shadow=true]{}
\parbox{#1\hsize}{%
      \[
        \textcolor{black} {#2}
      \]}
  \end{beamerboxesrounded}
}}

\AtBeginSection[] {
  \begin{frame}<beamer>\frametitle{Outline}
    \tableofcontents[currentsection,hideothersubsections]
  \end{frame}
}
%
%
%\AtBeginSubsection[] {
%  \begin{frame}<beamer>\frametitle{Outline}
%    \tableofcontents%[currentsection,currentsubsection]
%  \end{frame}
%}

%\usecolortheme[RGB={82,45,128}]{structure}
%\usecolortheme[RGB={162,80,22}]{structure}
\usecolortheme[RGB={128,0,0}]{structure}
\usetheme[secheader]{Boadilla}
%\usetheme[height=7mm]{Rochester}
%\usetheme{Copenhagen}
%\usetheme{Antibes}
%\usecolortheme{seahorse}
%\usecolortheme{crane}
%\usecolortheme{rose}
%\usefonttheme[onlylarge]{structurebold}
%\usefonttheme[onlymath]{serif}



\def\diag{{\rm diag}}


\def\E{\mathbb{E}}
\def\Prob{\mathbb{P}}
\def\argmin{{\rm argmin}}
\def\argmax{{\rm argmax}}
\def\Def{\stackrel{def}{=}}


\newtheorem{assumption}{Assumptions}
\newtheorem*{proposition}{Proposition}
\newtheorem*{remark}{Remark}



%\setbeamercolor{disc title}{bg=oiBlue!40!white!60,fg=blue}
\setbeamercolor{disc body}{bg= Regalia!20!white!80,fg= Regalia!80!black!90}

\setbeamercolor{clicker ungraded title}{bg=irishGreen!80!white!50,fg=irishGreen!30!black!90}
\setbeamercolor{clicker ungraded body}{bg=irishGreen!20!white!80,fg=irishGreen!30!black!90}

\setbeamercolor{clicker review title}{bg=gray!80!white!80,fg=oiBlue!80!black!90}
\setbeamercolor{clicker review body}{bg=gray!30!white!90,fg=oiBlue!80!black!90}

\setbeamercolor{code body}{bg=gray!20!white!80,fg=black}


% Custom commands
\newcommand{\degree}{\ensuremath{^\circ}}
\newcommand{\Note}[1]{
\rule{2.5cm}{0.25pt} \\ \textit{\scriptsize {\textcolor{rubineRed}{Note:} \textcolor{gray}{#1}}}}
\newcommand{\ct}[1]{
\vfill
{\tiny #1}}
\newcommand{\Remember}[1]{\textit{\scriptsize{\textcolor{orange}{Remember:} \textcolor{gray}{#1}}}}
\newcommand{\red}[1]{\textit{\textcolor{rubineRed}{#1}}}
\newcommand{\pink}[1]{\textit{\textcolor{rubineRed!90!white!50}{#1}}}
\newcommand{\green}[1]{\textit{\textcolor{irishGreen}{#1}}}
\newcommand{\webURL}[1]{\urlstyle{same}\textit{\textcolor{linkCol}{\url{#1}}} }
\newcommand{\webLink}[2]{\href{#1}{\textcolor{linkCol}{{#2}}}}
\newcommand{\mail}[1]{\href{mailto:#1}{\textit{\textcolor{linkCol}{#1}}}}
\newcommand{\hl}[1]{\textit{\textcolor{oiBlue}{#1}}}
\newcommand{\hlGr}[1]{\textit{\textcolor{lightGreen}{#1}}}
\newcommand{\mathhl}[1]{\textcolor{oiBlue}{\ensuremath{#1}}}
\newcommand{\ex}[1]{\textcolor{blue}{{{\small (#1)}}}}
\newcommand{\disc}[1]{
\begin{beamerboxesrounded}[shadow = true, lower = disc body, upper = disc title]{}
#1
\end{beamerboxesrounded}
}

\newcommand{\cl}[1]{
\begin{beamerboxesrounded}[shadow = true, lower = clicker ungraded body, upper = clicker ungraded title]{Question}
$\:$ \\
#1
\end{beamerboxesrounded}
}

\newcommand{\clR}[1]{
\begin{beamerboxesrounded}[shadow = true, lower = clicker review body, upper = clicker review title]{\red{Review question} }
$\:$ \\
#1
\end{beamerboxesrounded}
}

\newcommand{\formula}[2]{
\begin{beamerboxesrounded}[shadow = true, lower = white, upper = clicker review body]{#1}
#2
\end{beamerboxesrounded}
$\:$ \\
}

\newenvironment{twocol}[4]{
\begin{columns}[c]
\column{#1\textwidth}
#3
\column{#2\textwidth}
#4
\end{columns}
}


\newenvironment{slot}[2]{
\begin{array}{c}
\underline{#1} \\
#2
\end{array}
}

\newcommand{\pr}[1]{
\left( #1 \right)
}

\newcommand{\solnMult}[1]{
\item[] \vspace{-0.59cm}
\only<beamer| beamer:1>{\item #1}
\soln{\only<2->{\item \red{#1}}}
}

%\newcommand{\codechunk}[1]{
%\begin{beamerboxesrounded}[shadow = true, lower = code body]{}
%{\small #1}
%\end{beamerboxesrounded}
%}

% Change margin

\newenvironment{changemargin}[2]{%
\begin{list}{}{%
\setlength{\topsep}{0pt}%
\setlength{\leftmargin}{#1}%
\setlength{\rightmargin}{#2}%
\setlength{\listparindent}{\parindent}%
\setlength{\itemindent}{\parindent}%
\setlength{\parsep}{\parskip}%
}%
\item[]}{\end{list}}

% Footnote

\long\def\symbolfootnote[#1]#2{\begingroup%
\def\thefootnote{\fnsymbol{footnote}}\footnote[#1]{#2}\endgroup}

% Commands from the book
\newenvironment{data}[1]{\texttt{#1}}{}
\newenvironment{var}[1]{\texttt{#1}}{}
\newenvironment{resp}[1]{\texttt{#1}}{}






%%%%%%%%%%%%%%%%%%%%%%%%%%%%%%%%%%%%%%%%%%%%%%%%%%%%%%%%%%%%%%%%%%%%%%%%%%%%%%%%%%%%%%%%%%%%%%%

\title[Chapter 4 part 1]{Chapter 4 part 1}
\subtitle{Discrete Random Variables}

%%%%%%%%%%%%%%%%%%%%%%%%%%%%%%%%%%%%%%%%%%%%%%%%%%%%%%%%%%%%%%%%%%%%%%%%%%%%%%%%%%%%%%%%%%%%%%%


\author[Jingchen (Monika) Hu] % (optional, use only with lots of authors)
{Jingchen (Monika) Hu}
% - Give the names in the same order as the appear in the paper.
% - Use the \inst{?} command only if the authors have different
%   affiliation.

\institute[Vassar] % (optional, but mostly needed)
{Vassar College}
% - Use the \inst command only if there are several affiliations.
% - Keep it simple, no one is interested in your street address.

\date[MATH 241] % (optional, should be abbreviation of conference name)
{MATH 241}
% - Either use conference name or its abbreviation.
% - Not really informative to the audience, more for people (including
%   yourself) who are reading the slides online

\subject{MATH 241}
% This is only inserted into the PDF information catalog. Can be left
% out.



% If you wish to uncover everything in a step-wise fashion, uncomment
% the following command:

%\beamerdefaultoverlayspecification{<+->}


\begin{document}




%%%%%%%%%%%%%%%%%%%%%

% Title Page

\begin{frame}%[plain]
\titlepage
\end{frame}

%%%%%%%%%%%%%%%%%%%%%
%\addtocounter{framenumber}{-1}
%
%\begin{frame}\frametitle{Annoucement}
%
%\begin{itemize}
%\item HW3: \red{due now!}
%\item HW4: \red{due Tuesday, Sept 30th}
%\end{itemize}
%
%
%\end{frame}
%
%


%%%%%%%%%%%%%%%%%%%%%
\begin{frame}{Outline}
%\tableofcontents[hideallsubsections,pausections]
\tableofcontents[hideallsubsections]
\end{frame}



%%%%%%%%%%%%%%%%%%%%%%%%%%%%%%%%%%%%%%%%%%
\section{Discrete random variables}
%%%%%%%%%%%%%%%%%%%%%%%%%%%%%%%%%%%%%%%%%%
\begin{frame}
\frametitle{Random Variables}

\begin{defn}
\hl{Random Variable} $X$ is a real-valued  function on the sample space $S$.
\end{defn}


\begin{itemize}
\item Random variable is a number associated with a random experiment.
\item Random variables are in essence a fancy way of describing an event.

\end{itemize}
\end{frame}



%%%%%%%%%%%%%%%%%%%%%%%%%%%%%%%%%%%%%%%%%%
\begin{frame}

\cl{Two fair six-sided dice are rolled. Let the random variable $X$ denote the product of the 2 dice. What are possible values of $X$ and their associated probabilities? Just give a few examples and you do not need to calculate the associated probability for all possible values.
}

\pause
\begin{align*}
P(X = 1) & = P(1,  1) = 1 / 6 \times 1 / 6\\
P(X = 2) & = P(1,  2) + P(2, 1) = 1 / 6 \times 1 / 6 \times 2\\
P(X = 3) & = P(1,  3) + P(3, 1) = 1 / 6 \times 1 / 6 \times 2\\
 \cdots & \cdots \\
P(X = 36) & = P(6,  6) = 1 / 6 \times 1 / 6
\end{align*}

\end{frame}


%%%%%%%%%%%%%%%%%%%%%%%%%%%%%%%%%%%%%%%%%%
\begin{frame}\frametitle{Cumulative distribution function (cdf)}
\begin{defn}
For a random variable $X$, the function $F$ defined by
\[ F(x) = P(X \leq x), \quad  -\infty < x < \infty \]
is called the \hl{cumulative distribution function} of $X$.
\end{defn}

\pause
Note that
\begin{itemize}
\item Capital $X$: random variable
\item Little $x$: a real-valued number
\item $\leq$: smaller than or equal to
\end{itemize}

\end{frame}

%%%%%%%%%%%%%%%%%%%%%%%%%%%%%%%%%%%%%%%%%%
\begin{frame}\frametitle{Properties of the cdf (Ch 4.10)}

\begin{itemize}
\item $P(a < X \leq b) = F(b) - F(a)$, for all $a<b$

\item $P(X < b)$ does not necessary equal to $P(X\leq b)$.
\end{itemize}

\begin{enumerate}
\item $F(x)$ is a non-decreasing function; i.e., if $x_1 < x_2$, then
\[F(x_1) \leq F(x_2)\]

\vspace{-0.5cm}
\item \[ \lim_{x \rightarrow \infty} F(x) = 1 \Longleftrightarrow P(X \leq \infty) = 1\]

\vspace{-0.5cm}
\item \[ \lim_{x \rightarrow -\infty} F(x) = 0 \Longleftrightarrow P(X > -\infty) = 1\]

\vspace{-0.1cm}
\item $F(x)$ is right continuous; i.e., for any decreasing sequence $\{x_n: n = 1, 2, \ldots \}$ that converges to $x$,
\[ \lim_{n\rightarrow \infty}F(x_n) = F(x)\Longrightarrow \lim_{n\rightarrow \infty}P(X \leq x + \frac{1}{n}) = P(X \leq x)\]
\end{enumerate}


\end{frame}


%%%%%%%%%%%%%%%%%%%%%%%%%%%%%%%%%%%%%%%%%%
\begin{frame}

\cl{Two fair six-sided dice are rolled. Let the random variable $X$ denote the product of the 2 dice. Find the probability of: (a) $P(X \leq 2)$, (b) $P(X \leq 35)$.
}

\pause
Part (a)
\begin{align*}
P(X \leq 2) &= P(X = 1) + P(X = 2) \\
&=  P(1, 1) + P(1, 2) + P(2, 1) = 1 / 6 \times 1 / 6 \times 3
\end{align*}

\pause
Part (b)
\begin{align*}
P(X \leq 35) & = P(X = 1) + P(X = 2) + \cdots + P(X = 35)\\
&= 1 - P(X = 36) \\
 &= 1 - P(6, 6) \\
&=  1 - 1 / 6 \times 1 / 6 
\end{align*}
\end{frame}




%%%%%%%%%%%%%%%%%%%%%%%%%%%%%%%%%%%%%%%%%%
\begin{frame}\frametitle{Discrete random variables}
\begin{defn}
\begin{dinglist}{\DingListSymbolA}
\item A random variable $X$ that can take on at most a countable number of possible values is a
\hl{discrete random variable}.
\item For a discrete random variable $X$, we define the \hl{probability mass function} (pmf) by
\[ p(x) = P(X = x) \]
\end{dinglist}
\end{defn}
\vspace{-0.1cm}
Note: pmf is also written as $f(x)$.


\end{frame}


%%%%%%%%%%%%%%%%%%%%%%%%%%%%%%%%%%%%%%%%%%
\begin{frame}\frametitle{pmf and cdf}
\begin{itemize}
\item For a discrete random variable $X$, there exists a countable sequence $x_1, x_2, \ldots$, such that
\begin{align*}
p(x_i) > 0 & \quad\text{ for } i = 1, 2, \ldots\\
p(x) = 0 & \quad\text{ for all other values of }x
\end{align*}

and
\[ \sum_{i=1}^\infty p(x_i) = 1 \]


\item Relationship between pmf and cdf (for discrete random variable)
\[ F(a) = \sum_{\text{all } x \leq a} p(x)\]

\item If we know pmf, we can compute cdf. And vice versa.

\end{itemize}

\end{frame}


%%%%%%%%%%%%%%%%%%%%%%%%%%%%%%%%%%%%%%%%%%
\begin{frame}

\cl{Three fair coins are tossed. Let the random variable $X$ denote the number of heads. Write out the pmf and cdf of $X$.
}

\pause
For pmf:
\begin{align*}
P(X = 0) &= P(TTT)  = (1 / 2)^3 = 1 / 8 \\
P(X = 1) &=  P(HTT) + P(THT) + P(TTH) = (1 / 2)^3 \times 3 = 3 / 8\\
P(X = 2) &= P(HHT) + P(HTH) + P(THH) = (1 / 2)^3 \times 3 = 3 / 8\\
P(X = 3) &= P(HHH) = (1 / 2)^3 = 1 / 8 
\end{align*}

\pause
For pmf:
\begin{align*}
F(a)= 
\begin{cases}
    1 / 8,& 0 \leq a < 1\\
    1 / 2,& 1 \leq a < 2 \\
    7 / 8,& 2 \leq a < 3 \\
    1, & 3 \leq a
\end{cases}
\end{align*}
\end{frame}


\begin{frame}

\cl{Suppose that the distribution function of $X$ is given by
\begin{align*}
F(b)= 
\begin{cases}
    0,& b < 0\\
    \frac{b}{4},& 0 \leq b < 1 \\
    \frac{1}{2} + \frac{b - 1}{4},& 1 \leq b < 2 \\
    \frac{11}{12}, & 2 \leq b < 3 \\
    1, & 3 \leq b
\end{cases}
\end{align*}
Find $P(X = i), i = 1, 2, 3$.
}

\pause

\begin{itemize}
\item $P(X = 1) = P(X \leq 1) - P(X < 1) = \frac{1}{2} - \frac{1}{4} = \frac{1}{4}$\\ \pause
\item $P(X = 2) = P(X \leq 2) - P(X < 2) = \frac{11}{12} - (\frac{1}{2} + \frac{2 - 1}{4}) = \frac{11}{12} - \frac{3}{4} = \frac{1}{6}$ \\ \pause
\item $P(X = 3) = P(X \leq 3) - P(X < 3) = 1 - \frac{11}{12} = \frac{1}{12}$
\end{itemize}

\end{frame}

%%%%%%%%%%%%%%%%%%%%%%%%%%%%%%%%%%%%%%%%%%
\begin{frame}\frametitle{Recap}

Random Variable
\begin{itemize}
\item Random Variable $X$ is a real-valued  function on the sample space $S$.
\[ X: S \longrightarrow \mathbb{R}\]

\item Cumulative distribution function (cdf)
\[ F_X(x) = P(X \leq x), \quad \text{for any } x \in \mathbb{R} \]
\end{itemize}

Discrete random variable
\begin{itemize}
\item can only take at most a countable number of possible values.
\item probability mass function (pmf)
\[ p_X(x) = P(X = x) \]
\end{itemize}

\vspace{-0.3cm}
\twocol{0.5}{0.5}
{
\begin{itemize}
\item $\sum_{i=1}^{\infty} p(x_i) = 1$
\end{itemize}
}
{
\begin{itemize}
\item $F(a) = \sum_{\text{all } x \leq a} p(x)$
\end{itemize}
}



\end{frame}

%%%%%%%%%%%%%%%%%%%%%%%%%%%%%%%%%%%%

%%%%%%%%%%%%%%%%%%%%%%
%\begin{frame}{Outline}
%%\tableofcontents[hideallsubsections,pausections]
%\tableofcontents[hideallsubsections]
%\end{frame}


%%%%%%%%%%%%%%%%%%%%%%%%%%%%%%%%%%%%%%%%%%
\section{Expectation}
%%%%%%%%%%%%%%%%%%%%%%%%%%%%%%%%%%%%%%%%%%
\begin{frame}
\frametitle{Expected value}
\begin{defn}
The \hl{expected value} (or \hl{mean}) of a \underline{discrete random variable} is defined as
\[E[X] = \sum_{x: p(x) > 0} x\cdot P(X = x)= \sum_{x: p(x) > 0} x p(x)\]
\end{defn}


\begin{itemize}
\item $E[X]$ is a weighted average of the possible values $x$ that $X$ can take on,
each value being weighted by the probability $p(x) $.
\end{itemize}

\vspace{0.5cm}
\twocol{0.2}{0.7}
{
\includegraphics[width = 0.5\textwidth]{figures/joke1}
\includegraphics[width = 0.5\textwidth]{figures/joke2}
}
{
When she told me I was average, \\
she was just being mean.
}


\end{frame}



%%%%%%%%%%%%%%%%%%%%%%%%%%%%%%%%%%%%%%%%%%
\begin{frame}
\cl{Toss a coin. Suppose the probability of a head is $p$. Let $X$ be
a 0-1 indicator random variable s.t.\
\[
X = \left\{
\begin{array}{ll}
1	& \text{if head is obtained}\\
0	& \text{otherwise}\\
\end{array}
\right.
\]
Compute $\mu = E[X]$.
}

%\invisible{
\pause
\[ E[X] = 1 \cdot p + 0 \cdot (1-p) = p\]
%}

\pause
\begin{itemize}
\item In general, for indicator variable $X = \delta_A$ or $\mathbf{1}_A$, denoted as
\[
X = \left\{
\begin{array}{ll}
1	& \text{if event } A \text{ occurs}\\
0	& \text{otherwise}\\
\end{array}
\right.
\]
\pause
The expected value of $X$ equals the probability that $A$ occurs.
\[ E[X] = 1 \cdot P(A) + 0 \cdot P(A^c) = P(A) \]
\end{itemize}


\end{frame}



%%%%%%%%%%%%%%%%%%%%%%%%%%%%%%%%%%%%%%%%%%
\begin{frame}
\cl{Let the random variable $X$ denote the GP a certain student will earn in this class. Suppose its pmf is
\[ p(0) = 0.05, \quad p(1) = 0.05, \quad p(2) = 0.3, \quad p(3) = 0.4\]
Calculate their expected GP $E[X]$.
}
\vspace{0.4cm}

\pause
Since the domain of $X$ (all possible values $X$ can take) is $\{0, 1, 2, 3, 4\}$,
first compute $p(4)$.
\[p(4) = 1 - p(0) - p(1) - p(2) - p(3) = 0.2\]
Then compute $E[X]$
\[ E[X] = 0 \times 0.05 + 1 \times 0.05 + 2 \times 0.3 + 3 \times 0.4 + 4 \times 0.2 = 2.65 \]


\end{frame}


%%%%%%%%%%%%%%%%%%%%%%%%%%%%%%%%%%%%%%%%%%
\begin{frame}\frametitle{Expectation of a function of a random variable}
\begin{itemize}
\item If $X$ is a discrete random variable, and $g$ is a real-valued function then the expectation (or expected value) of $Y = g(X)$ is
\[ E[g(X)] = \sum_{x: p_X(x) > 0} g(x) \cdot p_X(x) \]
\end{itemize}


\end{frame}

%%%%%%%%%%%%%%%%%%%%%%%%%%%%%%%%%%%%%%%%%%
\begin{frame}
\cl{Let $X$ denote a random variable that takes on any of the values -1, 0, and 2 with respective probability: $P(X = -1) = \frac{1}{5}, P(X = 0) = \frac{1}{5}, P(X = 2) = \frac{3}{5}$. Compute $E[X^3]$.}

\pause
Let $Y = X^3$, then
\begin{align*}
E[X^3] = E[Y] & = \sum_{\text{all }x} x^3 \cdot p_X(x)\\
& = (-1)^3 \times \frac{1}{5} + 0^3 \times \frac{1}{5} + 2^3 \times \frac{3}{5}\\
& = -\frac{1}{5} + 0 + \frac{24}{5} \\
& = \frac{24}{5}
\end{align*}


\end{frame}

%%%%%%%%%%%%%%%%%%%%%%%%%%%%%%%%%%%%

\begin{frame}
\frametitle{Properties of expected values}

If $a$ and $b$ are constants, then
\[ E[aX+b] = a E[X] + b \]


\begin{itemize}
\item Holds for all random variable $X$ (not necessary discrete random variable).
\item Special cases of linear transformation $ E[aX+b] = a E[X] + b $
\begin{itemize}
\item constant factor \[ E[aX] = a E[X] \]

\item constant \[ E[b] = b \]
\end{itemize}
\end{itemize}

\end{frame}

%%%%%%%%%%%%%%%%%%%%%%%%%%%%%%%%%%%%
\begin{frame}
\cl{Let $X$ denote a random variable that takes on any of the values -1, 0, and 2 with respective probability: $P(X = -1) = \frac{1}{5}, P(X = 0) = \frac{1}{5}, P(X = 2) = \frac{3}{5}$. Compute: (a) $E[2X^2]$, (b) $E[4X^2 - 1]$.}

\pause
Let $Y = X^2$, then
\begin{align*}
E[X^2] = E[Y] & = \sum_{\text{all }x} x^2 \cdot p_X(x)\\
& = (-1)^2 \times \frac{1}{5} + 0^2 \times \frac{1}{5} + 2^2 \times \frac{3}{5}\\
& = \frac{1}{5} + 0 + \frac{12}{5} \\
& = \frac{13}{5} 
\end{align*}

Part (a): $E[2X^2] = 2 E[X^2] = 2 E[X^2] = 26 / 5$\\
Part (b): $E[4X^2 - 1] = 4 E[X^2] - 1 = 52 / 5 - 1 =  47 / 5$
\end{frame}
%%%%%%%%%%%%%%%%%%%%%%%%%%%%%%%%%%%%


%%%%%%%%%%%%%%%%%%%%%%%%%%%%%%%%%%%%%%%%%%
\begin{frame}\frametitle{Recap}

Expectation $\mu$
\begin{itemize}
\item \textbf{For discrete random variable}: $E[X] = \sum_{\text{all } x} x \cdot p(x)$
\item \textbf{Functions}: $E[g(X)] = \sum_{\text{all }x} g(x)~p(x)$
\vspace{2mm}
\item \textbf{Indicators}: $E[\delta_A] = P(A)$ where $\delta_A$ is an indicator function
\item \textbf{Linear function}: $E[aX + b] = aE[X] + b$
\item \textbf{Constants}: $E[c] = c$ if $c$ is constant
\end{itemize}

For two random variable's $X$ and $Y$
\[E[X + Y] = E[X] + E[Y]\]

\end{frame}

%%%%%%%%%%%%%%%%%%%%%%%%%%%%%%%%%%%%

%%%%%%%%%%%%%%%%%%%%%%
%\begin{frame}{Outline}
%%\tableofcontents[hideallsubsections,pausections]
%\tableofcontents[hideallsubsections]
%\end{frame}


\end{document}
